\section{A reserve that survives adding either common output}
\label{reserve:section}

The linear rate $2I$ leaves a remainder that has a simple behavior
under mixing.  Write $S=F-2I$, as in
\eqref{scalar:surplus-definition}, and put
\begin{equation}
 R(v,e)=vS\!\left(L-\frac ev\right),\qquad v,e>0,
 \label{reserve:perspective}
\end{equation}
using the zero extension of $S$ at negative arguments.  The profile
\eqref{main:profile} is
\[
 \mathcal B_m(I)=2I+R(1-m^2,H(m)-I).
\]
Its construction can be read as an allocation of information between
two capacities.  Set $\alpha=1-m^2$ and
$c_m=H(m)-\alpha L\ge0$.  Then
\begin{equation}
 \mathcal B_m(I)=
 \inf_{\substack{u+v=I\\0\le u\le c_m\\0\le v<\alpha L}}
 \{2u+\alpha F(v/\alpha)\},\qquad 0\le I<H(m).
 \label{reserve:allocation}
\end{equation}
Indeed, $F'\ge2$ makes the optimum $u=\min\{I,c_m\}$.
The capacity $c_m$ is paid at the linear rate, while the remaining
capacity $\alpha L$ uses the exact coordinate rate.  This is an
identity for the profile, not a decomposition of an actual channel.
The finite-noise bound of Kramer and Saglam
\cite[Theorem~1.1]{KramerSaglam2026} reads, in the present
normalization,
\[
 I(T_\rho f)\le\rho^2c_m+\alpha\psi(\rho)
 \qquad(0\le\rho\le1)
\]
for a Boolean field $f$ on the uniform cube with $\E f=m$.
It uses the same coefficients $\alpha$ and $c_m$; here they
define the local production profile $\mathcal B_m$.

To determine the payment required by a join, consider a physical state
\[
 |m|+\delta<1,\quad\delta\ge0,\quad
 0<e\le\bar h:=\frac{H(m+\delta)+H(m-\delta)}2,
\]
and write
\[
 \alpha=1-m^2,\quad\beta=\alpha-\delta^2,
 \qquad J=H(m)-\bar h,\quad j=\bar h-e.
\]
Set $a=m+\delta$ and $b=m-\delta$.  Let $e_a,e_b>0$ satisfy
$e_a\le H(a)$, $e_b\le H(b)$, and $(e_a+e_b)/2=e$, and put
$I_a=H(a)-e_a$, $I_b=H(b)-e_b$.  Their average information is
$j$, and their average variance capacity is
$[(1-a^2)+(1-b^2)]/2=\beta$.  Since $R$ is a convex
perspective, the child profiles satisfy
\begin{equation}
 \frac{\mathcal B_a(I_a)+\mathcal B_b(I_b)}2
 \ge2j+R(\beta,e).
 \label{reserve:pooling}
\end{equation}
The parent information is $I=H(m)-e=J+j$.  Subtracting the
pooled bound from $\mathcal B_m(I)=2(J+j)+R(\alpha,e)$ gives
the required payment
\begin{equation}
 W(m,\delta,e)=2J+R(\alpha,e)-R(\beta,e).
 \label{reserve:payment}
\end{equation}
Thus
\[
 \mathcal B_m(I)-\frac{\mathcal B_a(I_a)+\mathcal B_b(I_b)}2
 \le W(m,\delta,e).
\]
It is enough to prove $W\le\mathcal K$.  The deterministic
capacity comparison pays the interior core; the next two lemmas
show that this payment survives adding either common output.

\subsection{A logarithmic bound for the corner geometry}

\begin{lemma}
\label{reserve:geometry}
For $-1<m<1$ and $0\le\delta<1-|m|$, put
\[
 A=1-m,\qquad
 \kappa=\frac12\log\frac1{1-\delta^2/(1+m)^2}.
\]
Then
\begin{equation}
 A^2\log\frac\alpha\beta\le3\kappa.
 \label{reserve:log-geometry}
\end{equation}
\end{lemma}

\begin{proof}
The case $\delta=0$ is immediate.  Otherwise set
\[
 h=\frac{1-m}{1+m}>0,\qquad
 t=\frac{\delta^2}{(1+m)^2}<\min\{1,h^2\},\qquad
 f(t)=-\log(1-t).
\]
The ratio $f(t/h)/f(t)$ decreases in $t$ when $h\ge1$ and
increases when $h\le1$.  Indeed, the ratio of the corresponding
positive integrands is $(1-t)/(h-t)$, whose derivative has sign
$1-h$; the same monotonicity passes to their integral ratio.

If $h\ge1$, the ratio is at most its limit $1/h$ at zero.
Thus $A^2f(t/h)/f(t)\le A^2/h=1-m^2\le1$.
If $0<h<1$, it suffices to bound the ratio at $t=h^2$.
The function
\[
 G(h)=(2-h)f(h)-(2+h)\log(1+h)
\]
has $G(0)=0$ and
\[
 G'(h)=2\left[\frac h{1-h^2}-V(h)\right]>0.
\]
The last sign follows by bounding the integral defining $V(h)$
by its endpoint integrand.  Consequently
\[
 \frac{f(h^2)}{f(h)}
 =1-\frac{\log(1+h)}{f(h)}\ge\frac{2h}{2+h},
\]
and
\[
 A^2\frac{f(h)}{f(h^2)}
 \le\frac{2h(2+h)}{(1+h)^2}
 =2-\frac2{(1+h)^2}\le\frac32.
\]
Since $f(t)=2\kappa$ and $f(t/h)=\log(\alpha/\beta)$,
both cases give \eqref{reserve:log-geometry}.
\end{proof}

\subsection{The common-output contraction}

\begin{lemma}[Adding either common output]
\label{reserve:corner}
For every physical state $(m,\delta,e)$ and $0<\tau\le1$,
\begin{equation}
 W(1-\tau+\tau m,\tau\delta,\tau e)
 \le\tau W(m,\delta,e).
 \label{reserve:plus-corner}
\end{equation}
The same inequality holds with parent mean $-1+\tau+\tau m$,
corresponding to the negative common output.  There is no sign
restriction on the original mean.
\end{lemma}

\begin{proof}
The positive-corner ray scales $A=1-m$, $\delta$, and $e$ by
the same factor.  Hold $\delta/A$ and $e/A$ fixed, and put
$\eta=1+\delta^2/A^2$.  Entropy differentiation gives
\begin{equation}
 A\frac{dJ}{dA}-J=\kappa.
 \label{reserve:entropy-euler}
\end{equation}
To check this identity, use
$H(m)-(1-m)V(m)=\log[2/(1+m)]$; the average of the two
child logarithms is the parent logarithm plus $\kappa$.

For $h>0$, define the entropy intercept
\[
 \gamma(h)=S(L-h)+hS'(L-h)
 =\begin{cases}
 \displaystyle\int_h^L\frac{P(u)}u\,du,&0<h<L,\\
 0,&h\ge L,
 \end{cases}
 \qquad P(u)=u^2F''(L-u).
\]
By the scalar curvature bound, $P\le M=4L^2/3$ and $2>3M$.
Thus $\gamma\ge0$, and
\begin{equation}
 0\le\gamma(e/\alpha)-\gamma(e/\beta)
 \le M\log\frac\alpha\beta.
 \label{reserve:intercept-bound}
\end{equation}
If an argument crosses the clipping point, the integral is simply
truncated at $L$.

Since $\alpha=A(2-A)$ and $\beta=A(2-\eta A)$, differentiation
of the perspectives in \eqref{reserve:payment} gives
\begin{equation}
 A\frac{dW}{dA}-W
 =2\kappa+A^2\{\eta\gamma(e/\beta)-\gamma(e/\alpha)\}.
 \label{reserve:euler}
\end{equation}
For example, with $e=Az$, one has
$R(\alpha,e)/A=(2-A)S(L-z/(2-A))$, whose derivative is
$-\gamma(e/\alpha)$.  The second perspective gives the factor
$\eta$.  The clipped remainder is $C^1$, because $S(0)=S'(0)=0$,
so these identities hold across the seam; no second derivative of
the zero extension is needed.

Now $\eta\ge1$, \eqref{reserve:intercept-bound}, and
Lemma~\ref{reserve:geometry} give the whole comparison:
\begin{equation}
 A\frac{dW}{dA}-W
 \ge2\kappa-MA^2\log\frac\alpha\beta
 \ge(2-3M)\kappa\ge0.
 \label{reserve:corner-margin}
\end{equation}
Therefore $W/A$ is nondecreasing.  Comparing $A$ and $\tau A$
proves \eqref{reserve:plus-corner}.  The intermediate states remain
physical by entropy concavity, and all child means stay interior for
$\tau>0$.  Finally $W$ is even in $m$, so reflection proves the
negative-corner statement.  For $\delta=0$, every displayed payment
and margin is zero.
\end{proof}

\subsection{From one core to the local inequality}

\begin{theorem}[Local joining for the linear reserve]
\label{reserve:joining}
Let $(A,B)$ have any joint law on $[-1,1]^2$ with
$e_a=\E H(A)>0$ and $e_b=\E H(B)>0$.  Put
\[
 a=\E A,\quad b=\E B,\quad m=\frac{a+b}{2},
 \quad e=\frac{e_a+e_b}{2},
\]
and $I=H(m)-e$, $I_a=H(a)-e_a$, $I_b=H(b)-e_b$.  Then
\begin{equation}
 \E c(A,B)\ge
 \mathcal B_m(I)-\frac{\mathcal B_a(I_a)+\mathcal B_b(I_b)}2.
 \label{reserve:local-bound}
\end{equation}
\end{theorem}

\begin{proof}
An infinite cost is immediate.  Otherwise positive child entropies
give interior means.  Swap the children if necessary so that
$\delta=(a-b)/2\ge0$; entropy concavity makes $(m,\delta,e)$
physical.

First pay the deterministic capacity state.  For an interior pair
$x,y$, let $m_0=(x+y)/2$, $\delta_0=|x-y|/2$,
$e_0=[H(x)+H(y)]/2$, and $J_0=H(m_0)-e_0$.
Write $\alpha_0=1-m_0^2$, $\beta_0=\alpha_0-\delta_0^2$.
The entropy lower bound gives $e_0\ge\beta_0L$ and
\[
 [L-e_0/\alpha_0]_+\le J_0/\alpha_0<L.
\]
The last strict inequality and the variance-refined one-bit bound
are Lemma~\ref{capacity:onebit}.  Hence
\begin{equation}
 \begin{aligned}
 W(m_0,\delta_0,e_0)
 &=2J_0+\alpha_0S(L-e_0/\alpha_0)\\
 &\le2J_0+\alpha_0S(J_0/\alpha_0)
 =\alpha_0F(J_0/\alpha_0)\le c(x,y).
 \end{aligned}
 \label{reserve:capacity}
\end{equation}

For $\delta>0$, Proposition~\ref{canonical:coverage} supplies
an exact optimizer with interior core $(x,y)$ of mass $k>0$
and equal-corner masses $\pi_+,\pi_-$, where
$k+\pi_++\pi_-=1$.  Its cost is $k c(x,y)$.
Starting from the core, add the positive corner with retention
factor $k/(k+\pi_+)$, then the negative corner with retention
factor $k+\pi_+$.  These two operations produce exactly the
canonical law.  Lemma~\ref{reserve:corner} and
\eqref{reserve:capacity} give
\begin{equation}
 W(m,\delta,e)
 \le(k+\pi_+)\frac{k}{k+\pi_+}W(m_0,\delta_0,e_0)
 \le k c(x,y)=\mathcal K(m,\delta,e).
 \label{reserve:canonical-payment}
\end{equation}
Both retention factors are positive; a missing corner gives a
factor one.  No orientation or separate case for the core is required.
For $\delta=0$, $W=\mathcal K=0$ directly.

Finally the actual pair law has cost at least
$\mathcal K(m,\delta,e)$.  Combining
\eqref{reserve:canonical-payment} with the original-child pooling
bound \eqref{reserve:pooling} gives
\[
 \E c(A,B)+\frac{\mathcal B_a(I_a)+\mathcal B_b(I_b)}2
 \ge2(J+j)+R(\alpha,e)=\mathcal B_m(I).
\]
All child profiles use their original individual entropies; only
the joining cost was minimized with their common average budget.
\end{proof}

The pooling step has a useful exact interpretation.  If
$e<\beta L$, its minimum over the entropy split is attained at
$e_a/(1-a^2)=e_b/(1-b^2)=e/\beta$; these entropies are
physical because $H(a)\ge(1-a^2)L$, and similarly for $b$.
If $e\ge\beta L$, allocate both child entropies in the intervals
$[(1-a^2)L,H(a)]$ and $[(1-b^2)L,H(b)]$ with the prescribed
average.  Both remainders vanish.  Thus the perspective in
\eqref{reserve:pooling} is the exact common-budget value of the
child reserve, rather than an additional approximation.
